\section{Relativistic Invariants}
\label{sec:relativistic_invariants}

The components of a four-vector depend on the inertial frame, but certain combinations of components do not. These frame-independent quantities are Lorentz invariants. They encode physical information that every inertial observer agrees upon.

The most common invariants are obtained by complete contraction of tensor indices. This section develops a practical method for identifying them and applies it to spacetime separation, four-velocity, four-momentum, plane-wave phase, and systems of particles.

\subsection{Lorentz scalars}

A Lorentz scalar \(S\) has the same value in every inertial frame:

\[
\boxed{
S'=S
}.
\]

A scalar may be a fundamental scalar field,

\[
\phi'(x')=\phi(x),
\]

or it may be constructed by contracting tensor indices.

Examples include

\[
A_{\mu}B^{\mu},
\]

\[
T^{\mu}{}_{\mu},
\]

and

\[
F_{\mu\nu}F^{\mu\nu}.
\]

A formula containing no free indices is a candidate scalar, but the objects being contracted must themselves transform tensorially.

\subsection{The invariant interval}

The displacement between two events is

\[
\Delta x^{\mu}
=
(c\Delta t,\Delta\mathbf{x}).
\]

Its squared norm is

\[
\boxed{
\Delta x_{\mu}\Delta x^{\mu}
=
\Delta s^{2}
=
c^{2}(\Delta t)^{2}
-
|\Delta\mathbf{x}|^{2}
}.
\]

Although \(\Delta t\) and \(\Delta\mathbf{x}\) change under a boost, their Minkowski combination does not.

The sign of the interval is therefore invariant. All inertial observers agree whether two events are timelike, null, or spacelike separated.

\subsection{Invariant causal order}

For timelike-separated events, there exists a rest frame in which the events occur at the same spatial position. Their order in time cannot be reversed by a proper orthochronous Lorentz transformation.

For spacelike-separated events, different inertial observers may disagree about which event occurred first. This does not violate causality because no subluminal signal can connect the events.

Thus the invariant interval separates physically meaningful causal ordering from frame-dependent coordinate ordering.

\subsection{Four-velocity norm}

For a massive particle,

\[
U^{\mu}
=
\gamma(c,\mathbf{v}).
\]

Its invariant norm is

\[
\boxed{
U_{\mu}U^{\mu}=c^{2}
}.
\]

Expanding,

\[
\begin{aligned}
U_{\mu}U^{\mu}
&=
\gamma^{2}(c^{2}-v^{2})\\
&=
c^{2}.
\end{aligned}
\]

The three-velocity is frame dependent, while the four-velocity norm is not.

\subsection{Four-momentum and invariant mass}

The four-momentum is

\[
p^{\mu}
=
\left(
\frac{E}{c},
\mathbf{p}
\right).
\]

Its squared norm is

\[
\begin{aligned}
p_{\mu}p^{\mu}
&=
\frac{E^{2}}{c^{2}}
-
|\mathbf{p}|^{2}.
\end{aligned}
\]

For a particle of invariant mass \(m\),

\[
\boxed{
p_{\mu}p^{\mu}=m^{2}c^{2}
}.
\]

Therefore

\[
\boxed{
E^{2}
=
|\mathbf{p}|^{2}c^{2}
+
m^{2}c^{4}
}.
\]

Energy and three-momentum change under a boost, but the combination on the left side of

\[
\frac{E^{2}}{c^{2}}-|\mathbf{p}|^{2}
=
m^{2}c^{2}
\]

is invariant.

\subsection{Rest energy}

In the particle's rest frame,

\[
\mathbf{p}=0.
\]

The energy-momentum relation gives

\[
E^{2}=m^{2}c^{4}.
\]

For positive energy,

\[
\boxed{
E_{0}=mc^{2}
}.
\]

This is the rest energy. It is not the total energy in every frame; rather, it is the energy measured in the frame where the particle's spatial momentum vanishes.

\subsection{Massless particles}

For a massless particle,

\[
m=0.
\]

Hence

\[
\boxed{
p_{\mu}p^{\mu}=0
}.
\]

The energy-momentum relation becomes

\[
\boxed{
E=c|\mathbf{p}|
}.
\]

The four-momentum is null. Its individual energy and momentum components remain frame dependent.

\subsection{Worked example: recovering the mass}

A particle has energy

\[
E=5.0\,\mathrm{GeV}
\]

and momentum magnitude

\[
|\mathbf{p}|=4.0\,\mathrm{GeV}/c.
\]

Using

\[
m^{2}c^{4}
=
E^{2}-|\mathbf{p}|^{2}c^{2},
\]

we obtain

\[
\begin{aligned}
m^{2}c^{4}
&=
(5.0\,\mathrm{GeV})^{2}
-
(4.0\,\mathrm{GeV})^{2}\\
&=
25\,\mathrm{GeV}^{2}
-
16\,\mathrm{GeV}^{2}\\
&=
9\,\mathrm{GeV}^{2}.
\end{aligned}
\]

Therefore

\[
\boxed{
mc^{2}=3.0\,\mathrm{GeV}
}.
\]

Every inertial observer obtains the same invariant mass even though the measured values of \(E\) and \(\mathbf{p}\) differ.

\subsection{Total four-momentum of a system}

For several particles, define the total four-momentum

\[
\boxed{
P^{\mu}
=
\sum_{a}p_{a}^{\mu}
}.
\]

Because sums of four-vectors are four-vectors, \(P^{\mu}\) transforms covariantly.

The invariant mass \(M\) of the complete system is defined by

\[
\boxed{
P_{\mu}P^{\mu}=M^{2}c^{2}
}.
\]

Equivalently,

\[
\boxed{
M^{2}c^{4}
=
E_{\mathrm{tot}}^{2}
-
|\mathbf{P}_{\mathrm{tot}}|^{2}c^{2}
}.
\]

The system invariant mass is not generally the sum of the individual masses. It also includes kinetic and interaction energy measured in the center-of-momentum frame.

\subsection{Worked example: two photons}

Consider two photons with equal energy \(E\), moving in opposite directions along the \(x\)-axis. Their four-momenta are

\[
p_{1}^{\mu}
=
\left(
\frac{E}{c},
\frac{E}{c},0,0
\right),
\]

and

\[
p_{2}^{\mu}
=
\left(
\frac{E}{c},
-\frac{E}{c},0,0
\right).
\]

Each photon is massless:

\[
p_{1}^{2}=p_{2}^{2}=0.
\]

The total four-momentum is

\[
P^{\mu}
=
p_{1}^{\mu}+p_{2}^{\mu}
=
\left(
\frac{2E}{c},0,0,0
\right).
\]

Therefore

\[
P_{\mu}P^{\mu}
=
\frac{4E^{2}}{c^{2}}.
\]

The invariant mass of the two-photon system is

\[
\boxed{
M
=
\frac{2E}{c^{2}}
}.
\]

A system composed entirely of massless particles can have nonzero invariant mass when its total four-momentum is timelike.

\subsection{Scalar products of four-momenta}

For two particles,

\[
p_{1}\cdot p_{2}
=
p_{1\mu}p_{2}^{\mu}
\]

is invariant.

Using

\[
(p_{1}+p_{2})^{2}
=
p_{1}^{2}+p_{2}^{2}+2p_{1}\cdot p_{2},
\]

one can relate the invariant mass of a system to pairwise scalar products.

These combinations are central in collision and scattering kinematics because they can be evaluated in whichever frame makes the calculation simplest.

\subsection{The center-of-momentum frame}

If the total four-momentum is timelike, there exists a frame in which

\[
\mathbf{P}_{\mathrm{tot}}=0.
\]

This is the center-of-momentum frame. In that frame,

\[
P^{\mu}
=
(Mc,\mathbf{0}),
\]

and

\[
E_{\mathrm{tot}}=Mc^{2}.
\]

The invariant mass is therefore the total energy of the system in its center-of-momentum frame divided by \(c^{2}\).

\subsection{Invariant plane-wave phase}

Define the wave four-vector

\[
k^{\mu}
=
\left(
\frac{\omega}{c},
\mathbf{k}
\right).
\]

The phase of a plane wave may be written as

\[
\boxed{
k_{\mu}x^{\mu}
=
\omega t-
\mathbf{k}\cdot\mathbf{x}
}.
\]

Because this is a contraction of a covector and vector, it is invariant:

\[
k'_{\mu}x'^{\mu}
=
k_{\mu}x^{\mu}.
\]

Different observers measure different frequency and wave vector, but they agree on the phase assigned to a physical event.

\subsection{Frequency and wave-vector transformation}

Since \(k^{\mu}\) is a four-vector, a boost in the \(x\)-direction gives

\[
\boxed{
\frac{\omega'}{c}
=
\gamma\left(
\frac{\omega}{c}-\beta k_{x}
\right)
}
\]

and

\[
\boxed{
k'_{x}
=
\gamma\left(
k_{x}-\beta\frac{\omega}{c}
\right)
}.
\]

Equivalently,

\[
\boxed{
\omega'
=
\gamma(\omega-vk_{x})
}.
\]

This is the relativistic origin of the Doppler effect.

\subsection{Invariant wave-vector norm}

The squared norm is

\[
k_{\mu}k^{\mu}
=
\frac{\omega^{2}}{c^{2}}
-
|\mathbf{k}|^{2}.
\]

For a massless wave satisfying

\[
\omega^{2}=c^{2}|\mathbf{k}|^{2},
\]

we have

\[
\boxed{
k_{\mu}k^{\mu}=0
}.
\]

For a massive relativistic field, the corresponding wave four-vector has a nonzero timelike norm related to the mass parameter.

\subsection{Invariant derivative combinations}

For a scalar field \(\phi\), the contraction

\[
\boxed{
\partial_{\mu}\phi\,
\partial^{\mu}\phi
}
\]

is invariant.

Expanding,

\[
\partial_{\mu}\phi\,
\partial^{\mu}\phi
=
\frac{1}{c^{2}}
\phi_{t}^{2}
-
|\nabla\phi|^{2}.
\]

This explains why the relativistic scalar-field Lagrangian uses the particular difference between time-derivative and spatial-gradient terms.

Likewise,

\[
\boxed{
\Box\phi
=
\partial_{\mu}\partial^{\mu}\phi
}
\]

is a scalar if \(\phi\) is a scalar.

\subsection{Invariant integration measure}

For a proper Lorentz transformation,

\[
\det\Lambda=1.
\]

The Jacobian of the coordinate change therefore satisfies

\[
\mathrm{d}^{4}x'
=
|\det\Lambda|\,\mathrm{d}^{4}x
=
\mathrm{d}^{4}x.
\]

Thus

\[
\boxed{
\mathrm{d}^{4}x
}
\]

is invariant under proper Lorentz transformations.

If the Lagrangian density \(\mathcal{L}\) is a Lorentz scalar, then the action

\[
S
=
\int\mathcal{L}\,\mathrm{d}^{4}x
\]

is Lorentz invariant.

This provides the structural principle used to construct relativistic field theories.

\subsection{Scalar, covariant, and invariant}

The terms are related but not identical.

\begin{itemize}
    \item A covariant equation is written using objects that transform consistently, so its form is valid in every inertial frame.
    \item A Lorentz scalar is a quantity whose numerical value is unchanged.
    \item An invariant may be a scalar or another frame-independent property, such as the sign of an interval.
\end{itemize}

For example,

\[
A^{\mu}=B^{\mu}
\]

is a covariant vector equation, but the individual components are not invariant. The contraction

\[
A_{\mu}A^{\mu}
\]

is an invariant scalar.

\subsection{A practical invariant-checking procedure}

To determine whether an expression is a Lorentz scalar:

\begin{enumerate}
    \item identify the transformation type of every factor;
    \item write all tensor indices explicitly;
    \item confirm that every index is fully contracted;
    \item verify that no free index remains;
    \item check that fixed background quantities do not secretly select a preferred frame;
    \item expand the expression in components as a structural check;
    \item when useful, calculate it in a second frame or use the Lorentz-matrix identity.
\end{enumerate}

Complete contraction is the main diagnostic, but it is valid only when the factors are genuine tensors.

\subsection{Common mistakes}

Typical errors include:

\begin{enumerate}
    \item calling one component of a four-vector invariant;
    \item assuming energy or three-momentum is separately frame independent;
    \item confusing invariant mass with the sum of constituent rest masses;
    \item forming a Euclidean square instead of a Minkowski contraction;
    \item treating an expression with a free index as a scalar;
    \item forgetting that a system of massless particles may have nonzero invariant mass;
    \item assuming that a coordinate formula is invariant merely because it looks the same in one frame.
\end{enumerate}

\subsection{What has been established}

Lorentz invariants are observer-independent quantities constructed most commonly through complete tensor contractions. Central examples are

\[
\Delta x_{\mu}\Delta x^{\mu},
\qquad
U_{\mu}U^{\mu},
\qquad
p_{\mu}p^{\mu},
\qquad
k_{\mu}x^{\mu},
\qquad
\partial_{\mu}\phi\partial^{\mu}\phi.
\]

The invariant norm of four-momentum gives

\[
E^{2}=|\mathbf{p}|^{2}c^{2}+m^{2}c^{4},
\]

while the invariant phase organizes the relativistic transformation of frequency and wave vector.

The final section consolidates the sign, index, and unit conventions used throughout the chapter and shows how formulas change when alternative conventions are adopted.